% Shared main-text experimental protocol; results are prospective.
\section{Experiments}
\label{sec:setup}

The forthcoming study uses the current OOLONG-Pairs profile. Runner, attributor,
and proposer use the same frozen \textbf{DeepSeek-V4-Flash} deployment and
decoding settings. Mining, proposal, validation, checkpointing, and final
evaluation are implemented; results for this study remain forthcoming.

\paragraph{Data and splits.}
\label{sec:envs-and-splits}
OOLONG-Pairs is reconstructed from the \texttt{trec\_coarse} portion of
OOLONG-synth~\citep{bertsch2025oolong}, using twenty predicates over each record
collection and a deterministic scorer for sets of \emph{user-ID pairs}.
Table~\ref{tab:study-settings} gives the configured allocation. Mining,
validation, and short test contain disjoint task instances on shared
short-context records. Short-test results therefore measure performance on
held-out task instances. The separate long evaluation assesses performance on
larger record collections, without claiming entirely unseen underlying question
texts. The current split is retained. Appendix~\ref{app:protocol} specifies the
dataset pin, reconstruction, and overlap.

\begin{table}[htbp]
\centering\small
\begin{tabular}{@{}lrrl@{}}
\toprule
Role & Tasks & Attempts/task & Use \\
\midrule
Short mining & 20 & 2 & Failure evidence \\
Short validation & 10 & 1 per harness & Joint promotion decision \\
Short test & 10 & 3 per condition & Held-out task instances \\
Long test & 40 & 3 per condition & Larger record collections \\
\bottomrule
\end{tabular}
\caption{Configured study. At most four admitted edits per round, three
optimization rounds, and patience three. Upstream short/long window labels are
8,192/262,144 tokens; these are not
measured lengths of the reconstructed prompts.}
\label{tab:study-settings}
\end{table}

\paragraph{Comparison and metrics.}
\label{sec:baselines}
The primary comparison is the exact saved starting harness \texttt{H0*R}
(evaluation condition \texttt{initial}) against the frozen edited harness
(\texttt{sh\_rlm}), on matching short and long task sets. Reference conditions
are the sparse \hzero{} (\texttt{b1}), the upstream-style \hzerostar{}
(\texttt{h0\_star}), and a pinned \emph{paper reconstruction} of
$\lambda$-RLM~\citep{roy2026ycombinator}. The evaluation condition list must be
selected explicitly: the runner's default reference grid omits \texttt{initial}.
Exact pair-set equality defines a pass. We additionally report recorded
pair-level precision, recall, F1, and available missing/extra-pair counts,
together with usage, cost, runtime, and completion coverage. These supplementary
metrics do not alter the promotion gate. Final evaluation uses three attempts
per task; individual model attempts have no controlled random seed.

\paragraph{Interpretation and planned extensions.}
Results will compare matching tasks with completed, failed, and missing attempt
counts reported separately. Repeated attempts and predicates sharing a context
are not independent corpus samples. A promotion under one-attempt validation
is a selection decision, not a significance test. Cross-environment evaluation
and weight-training comparisons remain future work: this profile materializes
only OOLONG-Pairs short and long splits. No completed transfer result or
long-context performance gain is claimed here.
